% Exported from Math++. Compile with pdflatex.
\documentclass[11pt,a4paper]{article}
\usepackage[T1]{fontenc}
\usepackage[utf8]{inputenc}
\usepackage{lmodern}
\usepackage{amsmath,amssymb,amsthm,mathtools}
\usepackage{graphicx,booktabs,xcolor}
\usepackage[hidelinks]{hyperref}
\usepackage[margin=25mm]{geometry}
\theoremstyle{plain}
\newtheorem{theorem}{Theorem}
\newtheorem{lemma}[theorem]{Lemma}
\newtheorem{corollary}[theorem]{Corollary}
\newtheorem{proposition}[theorem]{Proposition}
\theoremstyle{definition}
\newtheorem{definition}[theorem]{Definition}
\newtheorem{example}[theorem]{Example}
\theoremstyle{remark}
\newtheorem*{remark}{Remark}
\setlength{\parindent}{0pt}\setlength{\parskip}{0.6em}
\title{Statistics}
\date{}
\begin{document}
\maketitle
\section*{A table in the document}

A cart rolls down a ramp; its position $x$ is measured every second.

\begin{center}
\begin{tabular}{ll}
\toprule
$t$ (s) & $x$ (m) \\
\midrule
0 & 1 \\
1 & 3 \\
2 & 5 \\
3 & 7 \\
\bottomrule
\end{tabular}
\end{center}
\begin{equation*}
  \bar{x} = 4\,\mathrm{m}
\end{equation*}
\begin{equation*}
  \sum_i x_i = 16\,\mathrm{m}
\end{equation*}
\begin{equation*}
  \operatorname{median}(x) = 4\,\mathrm{m}
\end{equation*}
\begin{equation*}
  \operatorname{std}(x) = \frac{2 \sqrt{15}}{3}\,\mathrm{m}
\end{equation*}
\begin{equation*}
  \operatorname{quartiles}(x) = \left(2\,\mathrm{m}, 4\,\mathrm{m}, 6\,\mathrm{m}\right)
\end{equation*}
\begin{equation*}
  \max(x) = 7\,\mathrm{m}
\end{equation*}
A straight line fits exactly, so $R^2 = 1$:

\begin{equation*}
  p(u) \coloneqq  \operatorname{linreg}(t, x)
\end{equation*}
\begin{equation*}
  p(10) = 21
\end{equation*}
\section*{Data from a file}

\begin{center}
\begin{tabular}{rr}
\toprule
$h$ (cm) & $m$ (kg) \\
\midrule
150 & 50 \\
160 & 56 \\
170 & 63 \\
180 & 70 \\
190 & 76 \\
\bottomrule
\end{tabular}
\end{center}
\begin{equation*}
  \bar{h} = 170\,\mathrm{cm}
\end{equation*}
\begin{equation*}
  \operatorname{var}(m) = 109\,\mathrm{kg^{2}}
\end{equation*}
\begin{equation*}
  \operatorname{pvar}(m) = \frac{436}{5}\,\mathrm{kg^{2}}
\end{equation*}
\begin{equation*}
  \operatorname{corr}(h, m) \approx 0.9995411791
\end{equation*}
\begin{equation*}
  \operatorname{ci}(m, 0.95) \approx [50.03664728 kg, 75.96335272 kg]
\end{equation*}
\section*{Distributions}

\begin{equation*}
  \operatorname{normalcdf}(1.96) \approx 0.9750021049
\end{equation*}
\begin{equation*}
  \operatorname{invnorm}(0.975) \approx 1.959963985
\end{equation*}
\begin{equation*}
  \operatorname{binompdf}(10, 0.5, 5) = \frac{63}{256}
\end{equation*}
\begin{equation*}
  \operatorname{tcdf}(2.228, 10) \approx 0.9749941141
\end{equation*}
\end{document}
